
% ============================================================================
\clearpage
\section*{v4.10.0-r63 Source-Only: The Koide Point as Fixed Point of the Frame Exchange, and What Remains Set or Fitted (k331--k339, H148--H156)}
\addcontentsline{toc}{section}{v4.10.0-r63 Source-Only: The Koide Point as Fixed Point of the Frame Exchange, and What Remains Set or Fitted (k331--k339, H148--H156)}
\label{sec:v4100-r63-koide-frame-exchange}
% ============================================================================

\paragraph{Purpose.}
Layer r62.8 found that the copy structure alone gives no masses, and r62.9 left a free mixing angle \(\theta\) between the \(HS(0)\) unit and
the copy direction.  On 25 and 26 November 2026 nine tests followed (H148--H156, scripts k331--k339), each fixed in an advance note before the
computation.  They ask what fixes \(\theta=\pi/4\) (Koide) and the Brannen phase \(\delta=2/9\), and trace the open choice step by step.  This
layer collects them with their status as recorded, on the author's decision of 26 November 2026.  It adds no new computation.  The outcome in one
sentence: \(K=\tfrac23\) follows from one postulate, the equal standing of the matter and the tachyon frame (exchange \(J\)), together with a set
scale; the postulate is not derived from the algebra, and \(\delta=2/9\) remains a fit.  The layer keeps apart throughout what is derived, what
is set and what is fitted (table in the balance below).  Signs are not set: where a sign matters, both cases are computed.  Earlier layers are
not rewritten.

\noindent\textbf{Source-only status.}
This v4.10.0-r63 layer is source-only.  It does not build a PDF, does not assign a new DOI,
and does not authorize a public release.  The v5.4.0 Version DOI
\href{https://doi.org/10.5281/zenodo.24500009}{10.5281/zenodo.24500009} and the Concept DOI
\href{https://doi.org/10.5281/zenodo.19210728}{10.5281/zenodo.19210728} are unchanged by it.

\subsection*{v4.10.0-r63 setting}
\label{subsec:v4100-r63-setting}

The state is \(\Psi=\cos\theta\,1+\sin\theta\,U\) with \(U\) a unit in the copy plane at azimuth \(\delta\); \(\theta\) is the polar angle between
the \(HS(0)\) unit and the copy plane, \(\delta\) the angle within the plane, and \(U'\) denotes the second direction of the plane.  The
\emph{amplitude readout} is that of r59.9: \(\sqrt{m_k}\propto c_k=\bigl(\cos\theta+\sqrt2\,\sin\theta\cos(\delta-2\pi k/3)\bigr)/\sqrt3\) with signed amplitudes, and
\(K=\sum m_k/(\sum\sqrt{m_k})^2\).  This readout is a setting.  In the potential, \(A=\varphi^2\) is the norm of the \(HS(0)\) component
(\(\varphi=\operatorname{Re}\Psi\)), \(B=|y|^2\) that of the copy direction with \(n_B\) components, and
\[
V=\mu_A^2A+\mu_B^2B+\lambda_AA^2+\kappa AB+\lambda_BB^2 .
\]
``Scale free'' means \(\mu_A^2=\mu_B^2=0\) with the scale arising by condensation at a Gildener--Weinberg point of the one-loop flow; ``scale
set'' means \(\mu^2\ne0\) as an input.  \(W\) is the diagonal element with entries \(c_k\) and \(N(W)=c_0c_1c_2\) its cubic norm.  Measured values
used: \(m_\tau=1776.93\pm0.09\) MeV and the pole masses of \(e\), \(\mu\).  Tools: Python (sympy, scipy, explicit octonions), Wolfram Language
independently for every part; all common quantities agree.  Qiskit was not used (no circuit).

\subsection*{v4.10.0-r63 results}
\label{subsec:v4100-r63-results}

\paragraph{Theoremlet v4.10.0-r63.1 (why the Gram values miss, and \(K=\tfrac23\) exactly at \(\theta=\pi/4\); H148, k331; partly not blind).}
The Gram readout of r62.8 compares fixed states of the three copies; their overlap is fixed by the \(120^\circ\) position that
\(\operatorname{Im}z^3=0\) forces (\(|a|^2=\tfrac1{16}\) for leptons, \(\tfrac14\) for colour), whereas \(K=\tfrac23\) needs
\(|a|^2=\tfrac12=\cos^2(\pi/4)\).  No choice of copies and no frame mask changes this (H148a confirmed).  In the amplitude readout
\[
K=\frac{1}{3\cos^2\theta}
\]
(Wolfram exact), independent of \(\delta\), so \(K=\tfrac23\) exactly at \(\theta=\pi/4\) (H148c confirmed).  With \(\delta=2/9\) inserted the
ratios are \(m_\mu/m_e=206.770\) (measured \(206.768\)) and \(m_\tau/m_\mu=16.8180\) (\(16.8177\)).  The complex charge state
\((e_0+ie_7)/\sqrt2\) gives \(K\approx0.34\) for every \(\delta\) (H148d); the Koide state must be real and hence fixed under time reversal.  A
Gram variant with \(\theta^*=26.2242^\circ\) meets \(K=\tfrac23\) with two equal masses; \(\theta^*\) is not a project quantity (H148b).
\emph{Derived:} \(K=\tfrac23\Leftrightarrow\theta=\pi/4\).  \emph{Set or fitted at this stage:} \(\theta=\pi/4\) itself, \(\delta=2/9\) (fitted to
\(m_\mu/m_e\)), the scale, the readout.

\paragraph{Theoremlet v4.10.0-r63.2 (\(\theta=\pi/4\) is the fixed point of the frame exchange; \(\delta\) is fixed by no term without fitting; H149, k332; part a not blind).}
\(\operatorname{Re}\Psi^2=\cos2\theta=x_0^2-|\vec x|^2\).  The exchange \(J:\psi\mapsto U\bar\psi\) maps \(\theta\) to \(\pi/2-\theta\), \(J^2=1\),
and its only fixed point in \([0,\pi)\) is \(\theta=\pi/4\), where \(\Psi\) is lightlike (\(\operatorname{Re}\Psi^2=0\)); verified in the octonions
explicitly.  By k297 \(J\) is the exchange of matter frame and tachyon frame (\(\mu\mapsto1/\mu\)).  For a quartic potential
\(\partial_\theta V|_{\pi/4}=-(2a-2c+\mu_0^2-\mu_1^2)/2\): exactly one linear condition, the exchange-odd part.  With the exchange it vanishes, and
the curvature is \(2(2a-b)\) at \(\pi/4\) and \(-2(2a-b)\) at \(\theta=0\) (\(K=\tfrac13\), three equal masses).  Minimum or maximum is decided by
the sign of \(2a-b\) alone; this sign is not set.  The HC breaks the exchange and gives \(K=(1+\beta_{\rm sym}^4)/3=0.364\) in the matter frame
(H149c).  For any \(V(N,\operatorname{Re})\) up to degree 4, \(\partial_\delta V\equiv0\); \(\delta\) enters only through
\(N(W)=(2\cos3\delta-\sqrt2)/(24\sqrt3)\) at \(\theta=\pi/4\).  Among 966 minima of \(V=\sum\alpha_kN^k\) with small integer coefficients, the 126
with three different masses all lie at \(\delta=\pi/12\) (\(N(W)=0\), massless electron, the octant of H101); none at \(2/9\) (H149d).  No project
quantity lies within \(1\%\) of \(\delta_{\rm fit}=0.222222\) (H149e).  Not evaluated: CT-002 (\(\delta=\pi/12-\gamma_M/6\)) lies \(5\cdot10^{-5}\)
off; a numerical coincidence \(\delta=\arccos(\pi/4)/3\) noticed before the note lies \(42\sigma\) off and does not hold.  H149a--e confirmed.

\paragraph{Theoremlet v4.10.0-r63.3 (the sign of \(2a-b=2\lambda-\kappa\) does not follow from the field equation alone; H150, k333; parts a, f not blind).}
With the potential of the full norm (or a gauged \(U(1)\)) \(2\lambda-\kappa=0\) to every order and \(\theta\) is flat.  \emph{Scale set:} the
Koide point is the minimum for the whole window \(-2\lambda<\kappa<2\lambda\); with a hard cutoff and unequal component numbers (1 against 2) the
mass term becomes exchange odd and \(\theta\) moves to \(\cos2\theta=c_L/f^2\) (this part of H150b \emph{refuted}).  \emph{Scale free:} a
condensate at the Koide point arises only for \(\kappa<0\) (class (iv); 190 of 190 for \(n=1\) and \(n=2\)), and then \(2a-b>0\) without a
further choice; for \(0<\kappa<2\lambda\) no scale arises.  H150c is \emph{refuted by its wording} in the part \(r>2\).  With the actual component
numbers (1 against 2) the flow shifts the condensation point to \(K\) between \(0.6667\) and \(0.6778\), median \(0.6719\) (H150d); with all of
\(\operatorname{Im}\mathbb O\) (1 against 7) the median is \(0.697\).  In the HC form \(\eta(B-\mu^2A)^2\) boundedness gives \(\kappa\le0\) and
\(2a-b>0\), but the minimum lies at \(\tan\theta=\mu\) (H150f).  So the sign follows from the field equation only if two things are required:
scale free, and the scale shall arise.  Deferred idea I4 (condensation as the source of mass) is met and was named to the author; it is not
implemented.

\paragraph{Theoremlet v4.10.0-r63.4 (the scale-free shift is start dependent and unrelated to any setting; H151--H153, k334--k336).}
(a)~The flow is invariant under rotations in the copy plane and does not generate \(\delta\) (H151a).  The shift of \(K\) is not a fixed number:
between \(0\) and \(0.011\) over the physical range, with a fixed value only in the limit \(\kappa_0\to0^-\),
\(K^*=\tfrac13+\sqrt{10}/9=0.684698\) for \(n_B=2\); no fixed ray with \(\kappa<0\) exists (H151b partly).  48 comparisons with functions of
\(2/9\) give no hit, 540 with all 15 settings give 2 against a chance expectation \(0.58\) (H151c, d).  Form S leaves \(K\) invariant and a boost
keeps the light cone fixed, so the frame difference cannot restore \(\tfrac23\) (H151e).
(b)~The complete scale-free ansatz with all \(\mathbb Z_3\rtimes\mathbb Z_2\)-invariant quartics (new only \(\xi\varphi C_3\),
\(C_3=y_1^3-3y_1y_2^2\)) and the back-reaction of \(HS(3)\): over 144 starting points \(K\) ranges from \(0.6615\) to \(0.8182\), and \(\delta\)
always lands on \(0\) or \(\pi/3\) (two equal masses); three different masses in 0 of 144 cases (H152a--e confirmed).
(c)~A normalisation of \(A\) and \(B\) (per component, by the trace \(\tfrac13:\tfrac23\), over three copies) is a reparametrisation (zero
residue): it selects only the starting ray (H153a, b).  Only restricting the loop to the subalgebra \(\{1,U\}\) (\(n_A=n_B=1\)) gives
\(K=\tfrac23\) and \(\theta=\pi/4\) exactly for all starts (H153c); that is a setting about the degrees of freedom.

\paragraph{Theoremlet v4.10.0-r63.5 (the loop over \(\{1,U\}\) is not forced; the deviation from \(\tfrac23\) is measurement noise; H154, k337; parts a, g, h not blind).}
The second plane direction contributes \(\beta(1,2)-\beta(1,1)=[2\kappa^2,\,8\kappa\lambda_B,\,8\lambda_B^2]\), nonzero in the Koide range.  The
only generator that leaves \(\operatorname{span}\{1,U\}\) invariant is the \(U(1)\) of \(\exp(\alpha U)\), under which \(U'\) is neutral; it
preserves \(V\) only for the full norm.  At the Gildener--Weinberg point \(U'\) is a Goldstone direction and \(K\) is unchanged; gauging the plane
rotation leaves \(K\) start dependent (\(0.634\) to \(1.922\) over 40 starts).  No symmetry removes a field from the loop: the algebra does not
force \(n_B=1\).  \emph{Measured value:}
\[
K=0.6666645\pm0.0000051,
\]
\(0.43\sigma\) below \(\tfrac23\), the error coming entirely from \(m_\tau\).  The value \(0.666661\pm0.000007\) used in H152 and H153 rested on
the older \(m_\tau=1776.86\pm0.12\) MeV; this is corrected here, and no decision changes.  One-loop QED running to a common scale shifts \(K\) by
\(+1.1\cdot10^{-3}\) (about \(225\sigma\)): Koide holds for pole masses, not for running masses at a common scale.  Form S leaves \(K\) exactly
unchanged and \(\partial_\delta K\equiv0\).  \emph{Conditional prediction:} if \(K=\tfrac23\) holds exactly for pole masses, then
\(m_\tau=1776.969\) MeV, \(0.039\) MeV above the present mean.  H154a--h confirmed.

\paragraph{Theoremlet v4.10.0-r63.6 (with a set scale the exchange symmetry holds \(K=\tfrac23\); H155, k338; part a not blind).}
For \(\mu_A^2=\mu_B^2=\mu^2<0\) and \(\lambda_A=\lambda_B=\lambda\) the minimum is \(A=B=-\mu^2/(2\lambda+\kappa)\), hence \(\theta=\pi/4\) and
\(K=\tfrac23\) at tree level for \(-2\lambda<\kappa<2\lambda\), with mass squares \(0\) (\(U'\), Goldstone), \(-4\mu^2\) and
\(4\mu^2(\kappa-2\lambda)/(\kappa+2\lambda)\); \(U'\) need not be heavy (H155a).  Making \(U'\) heavy by the cubic norm \(\zeta N(W)\) shifts \(K\)
together with its mass (H155b).  Running with two plane directions: the first order in the logarithm cancels exactly between mass and quartic
running (symbolic zero residue), leaving an effect of order \((\lambda L/16\pi^2)^2\); H155c, which predicted that every window misses, is
\emph{refuted} in favour of Koide.  A hard cutoff or thresholds shift \(\theta\) at first order (r63.3 stands).  A sixth-order term
\(\zeta N+N^2/\Lambda^2\) allows interior \(\delta\) for \(R=\zeta\Lambda^2/v^3\in[-0.0282,\,0.164]\); \(\delta=2/9\) requires
\(R^*=-0.0075807\), one parameter, which is a fit.  One comparison with known scales (\(R=-\sum m_\ell/v\), \(0.89\%\) off) counts as a hit by the
wording of H155d but lies \(63\sigma\) off in \(\delta\) and is not independent of the target; it is no derivation.  \emph{Exploration, after the
fact:} with \(R^*\) the full potential keeps \(K=\tfrac23\) and gives \(\delta=2/9\), doubly degenerate with \(\delta=0.29713\).  The three lepton
masses are then reproduced from the set exchange symmetry, the scale and one fitted parameter: as many inputs as Koide plus Brannen phase.

\paragraph{Theoremlet v4.10.0-r63.7 (equal mass terms follow only from the exchange \(J\) as a postulate; H156, k339; parts a, d not blind).}
\(\mu_AA+\mu_BB=\tfrac12(\mu_A+\mu_B)N+\tfrac12(\mu_A-\mu_B)\operatorname{Re}\Psi^2\), and both \(N\) and \(\operatorname{Re}\Psi^2\) are invariant
under conjugation and under \(G_2\): the algebra allows unequal mass terms (H156a).  Among the candidate symmetries, \(SO(2)=\exp(\alpha U)\) is
too strong (\(\kappa=2\lambda\), no Koide condensate), \(\mathbb Z_4=\langle L_U\rangle\) too weak (it allows the exchange-odd quartic
\(\varphi y(y^2-\varphi^2)\)); exactly \(\mathbb Z_2=\langle J\rangle\) and \(D_4=\langle L_U,\text{conjugation}\rangle\) enforce the form of
r63.6: equal masses, equal \(\lambda\), \(\kappa\) free (H156b).  \(J\) is an isometry but neither an automorphism nor an antiautomorphism:
\(J(1)J(1)=-1\ne J(1)\), \(J(zw)=J(z)\bar w\) (H156d).  The algebra therefore does not supply \(J\); it has to be required as a symmetry of the
field theory.  The flow helps only weakly: for \(n_B=1\) the quartic direction is attracted to the exchange line for \(|\kappa|<6\lambda\), the
mass direction only for \(\kappa<0\); for \(n_B=2\) the flow sources the deviation itself (H156c partly).  Equal mass terms are thus equivalent to
the \emph{postulate of equal standing of the matter and the tachyon frame}.  The author adopted this postulate as a working assumption on
26 November 2026; it is a statement about the physics, not a result of the computation.  Limit: the invariants are computed in the plane
\(\{1,U\}\); in all of \(\mathbb O\) or \(\mathbb S\), \(L_U\) leads out of the three-field space (named only).

\subsection*{v4.10.0-r63 balance of settings}
\label{subsec:v4100-r63-balance}

\begin{center}
\small
\begin{tabular}{@{}>{\raggedright\arraybackslash}p{0.30\textwidth}>{\raggedright\arraybackslash}p{0.62\textwidth}@{}}
\toprule
item & status\\
\midrule
\(K=\tfrac23\Leftrightarrow\theta=\pi/4\) & derived (in the amplitude readout)\\
\(\theta=\pi/4\) is the fixed point of \(J\); stationary for every exchange-symmetric potential & derived\\
\(\theta=\pi/4\) realised & only with the postulate of equal standing of the frames (\(J\)); postulate adopted by the author, not derived\\
scale \(\mu^2\) & set (scale free: \(K\) start dependent, shifted)\\
sign of \(2\lambda-\kappa\), of \(\mu^2\), \(\zeta\), global orientation & not set; both cases computed; left to the field equation\\
\(\delta=2/9\) & fitted (one parameter \(R^*\), or CT-002 as placeholder)\\
amplitude readout, number of generations & set (inputs)\\
\(m_\tau=1776.969\) MeV & conditional prediction (if \(K=\tfrac23\) exactly for pole masses)\\
\bottomrule
\end{tabular}
\end{center}
The open choice moved from ``\(\theta\)'' (r63.1) to ``sign of \(2a-b\)'' (r63.2), to ``starting value of the flow'' (r63.3--r63.5), to ``equality
of the mass terms'' (r63.6) and ends at one postulate (r63.7).  The target \(\delta=2/9\) is itself a fitted input; the question of \(\theta\)
arises in the chosen decomposition (amplitude readout) and would not arise in this form in the Gram readout.

\subsection*{v4.10.0-r63 frames}
\label{subsec:v4100-r63-frames}

\begin{center}
\small
\begin{tabular}{@{}>{\raggedright\arraybackslash}p{0.34\textwidth}>{\raggedright\arraybackslash}p{0.22\textwidth}>{\raggedright\arraybackslash}p{0.36\textwidth}@{}}
\toprule
item & frame & comparison\\
\midrule
r63.1 (readouts, measured ratios) & M & form S for charged leptons: \(K\) unchanged, so the comparison with data is frame independent\\
r63.2, r63.7 (exchange \(J\)) & M \(\leftrightarrow\) T, declared map \(\theta\mapsto\pi/2-\theta\) (\(\mu\mapsto1/\mu\), k297) & fixed point \(\theta=\pi/4\); HC value \(0.364\) in M\\
r63.3--r63.6 (potential, flow) & M & boost keeps the light cone; no rapidity restores \(\tfrac23\) (H151e)\\
P & not used & ---\\
\bottomrule
\end{tabular}
\end{center}
\(J\) is the exchange of matter and tachyon frame and is not time reversal \(z\mapsto\bar z\); the two are kept apart (r51, r81).

\subsection*{v4.10.0-r63 register entries}
\label{subsec:v4100-r63-register}

\begin{center}
\small
\begin{tabular}{@{}>{\raggedright\arraybackslash}p{0.40\textwidth}>{\raggedright\arraybackslash}p{0.52\textwidth}@{}}
\toprule
statement & status after v4.10.0-r63\\
\midrule
H148a--d; H149a--e; H150a, d, e, f; H151a, c, d, e; H152a--e; H153a, c, d; H154a--h; H155a, b; H156a, b, d & confirmed, with the blindness recorded per part\\
H150b (cutoff part, 1 against 2), H150c (part \(r>2\)), H155c, H155d & refuted (H150c, H155d by wording; H155c in favour of Koide)\\
H151b, H156c & partly\\
H153b & confirmed except one sub-prediction (two normalisations give the same \(K\))\\
Koide \(K=\tfrac23\) & follows from the postulate \(J\) with a set scale; not derived from the algebra\\
postulate of equal standing of M and T frame & adopted by the author on 26 November 2026 as working assumption\\
\(\delta=2/9\) & fit; CT-002 remains a placeholder\\
\(m_\tau=1776.969\) MeV & conditional prediction, open\\
consequences of \(J\) for \(\delta\) (H157) & in progress elsewhere; not part of this layer\\
OP3, OP4 & open\\
\bottomrule
\end{tabular}
\end{center}

\subsection*{v4.10.0-r63 notes}
\label{subsec:v4100-r63-notes}

{\raggedright
\texttt{notes/\allowbreak research/\allowbreak hilbert\_\allowbreak shift\_\allowbreak felder\_\allowbreak 2026\_\allowbreak 11\_\allowbreak 24/\allowbreak generationszahl\_\allowbreak h138/}:
advance notes \texttt{Vorab\_\allowbreak Notiz\_\allowbreak H148\_\allowbreak *.md} to \texttt{Vorab\_\allowbreak Notiz\_\allowbreak H156\_\allowbreak *.md}, result notes
\texttt{Ergebnis\_\allowbreak H148\_\allowbreak *.md} to \texttt{Ergebnis\_\allowbreak H156\_\allowbreak *.md}, scripts \texttt{k331} to \texttt{k339} (each with a
Wolfram cross-check \texttt{b}) and their actual outputs \texttt{k3NN\_\allowbreak out.txt}.  The exchange as frame change (k297) is in
\texttt{feldgleichung/\allowbreak HC\_\allowbreak Tauschsymmetrie\_\allowbreak k297.md} of the same folder tree.
\par}

\subsection*{v4.10.0-r63 ledger and guard update}
\label{subsec:v4100-r63-ledger-guard}

\[
R_{\rm chk}^{v4100\text{-}r63}=(1,1,0)
\]
for (status recorded; computations checked in the cited scripts; no new computation in this layer; \(\delta\), the scale and the postulate \(J\) not derived).

The layer spends no guard:
\[
R_{\rm glob}^{(v4100\text{-}r63,+)}=(0,1,1,1,1),
\qquad
\Gamma_{\rm guard}^{v4100\text{-}r63}=(1,1,1,1,1),
\qquad
\beta_4=1/256\ \text{unspent}.
\]
The layer does not imply exact finite \(C5\) zero closure, exact all-readout terminality, \(C6/T2\), a global photon-clock/RH theorem,
or global all-readout closure.  OP1--OP5, C5 and C6/T2 remain open.
